\section{Notation and conventions}\label{sec:prelim}

\subsection{Dirichlet characters and \texorpdfstring{$L$}{L}-functions}
For $N\ge1$ a Dirichlet character modulo $N$ is a homomorphism $\chi:(\Z/N\Z)^\times\to\C^\times$,
extended by $0$ to $\Z/N\Z$ and viewed as a function on $\Z$. The trivial character modulo $N$ is
$\one_N$. For $M\mid N$ and $\psi$ modulo $M$ we write $\psi_{\uparrow N}$ for the induced character
modulo $N$. The conductor $\cond(\chi)$ is the least $M\mid N$ such that $\chi=\psi_{\uparrow N}$ for
some $\psi$ modulo $M$; this $\psi=\chi^*$ is unique, and $\chi$ is \emph{primitive} if
$\cond(\chi)=N$. The only primitive trivial character is $\one_1$. We write $\chi$ \emph{even} or
\emph{odd} according as $\chi(-1)=1$ or $-1$, and $\mathfrak a(\chi)\in\{0,1\}$ accordingly.

$L(s,\chi)$ denotes the meromorphic continuation of $\sum_n\chi(n)n^{-s}$ to $\C$, following
Mathlib's conventions: it is a total function $\C\to\C$, entire for $\chi\ne\one_N$, and equal to
$\zeta(s)$ for $\chi=\one_1$, with the value $\zeta(1)=\frac12(\gamma-\log 4\pi)$ at the pole.
Put $\Gamma_\R(s)=\pi^{-s/2}\Gamma(s/2)$, $G_\chi(s)=\Gamma_\R(s+\mathfrak a(\chi))$ and
$\Lambda(\chi,s)=G_\chi(s)L(s,\chi)$. For $\chi\ne\one$ let
$\ell(\chi)=L'(1,\chi)/L(1,\chi)$; $\gamma$ is Euler's constant and $\Lambda(n)$ the von Mangoldt
function. For an odd prime $q$, $(\frac{\cdot}{q})$ is the Legendre symbol.

\subsection{Primitive characters}
A \emph{primitive character} is a pair $\psi=(N_\psi,\chi_\psi)$ with $\chi_\psi$ primitive modulo
$N_\psi$. We write $\mathcal P$ for the set of primitive characters, $\psi(n)=\chi_\psi(n)$,
$L(s,\psi)=L(s,\chi_\psi)$, $\one=(1,\one_1)$ and $\psi^{-1}=(N_\psi,\chi_\psi^{-1})$. For any character
$\chi$ we write $\chi^{\rm prim}=(\cond\chi,\chi^*)\in\mathcal P$.

\subsection{Number fields and Galois cohomology}
For a number field $E$, an algebraic closure $\Omega$ and a prime $\ell$, $S_\ell$ is the set of places
above $\ell$ and $G_{E,S_\ell}$ the Galois group of the maximal extension of $E$ in $\Omega$ unramified
outside $S_\ell$ (archimedean places are ignored throughout). For $j\in\Z$ we write $\Ql/\Zl(-j)$ for the
twist by the $(-j)$-th power of the cyclotomic character and
\[
\Sha^1(G_{E,S_\ell},\Ql/\Zl(-j))=\ker\Big(H^1(G_{E,S_\ell},\Ql/\Zl(-j))\to\prod_{v\in S_\ell}H^1(E_v,\Ql/\Zl(-j))\Big),
\]
with $\Zl$-length $\mathrm{sh}_\ell(E,j)\in\N\cup\{\infty\}$. By Poitou--Tate duality this group is
dual to $\Sha^2(G_{E,S_\ell},\Zl(1+j))$, the $\ell$-part of an \'etale wild kernel
\cite{NeukirchSchmidtWingberg2008}.

\subsection{Unitary groups and theta lifts}
If $E/F$ is a CM extension of number fields, a hermitian space $X$ over $E$ is \emph{totally definite}
if it is positive definite at every real place of $F$; $\Uni(X)$ is its unitary group. For a pair of
hermitian and skew-hermitian spaces the Weil representation of the corresponding unitary dual pair, and
the associated global theta lift, are as in \cite{Howe1989,KudlaRallis1994}; the Siegel--Weil formula is
used in the form of \cite{Weil1965,KudlaRallis1988,Ichino2004}. The finer data used in
\cref{sec:families} are introduced there.

\subsection{Conventions}
A set $W\subset\R$ \emph{has windows of exponent $\lambda>1$} if for every $y$ there is $Y\ge y$ with
$[Y,Y^\lambda]\subset W$. ``Eventually in $x$'' means ``for all $x$ larger than some $x_0$''. Constants
$C,c,\dots$ may change from line to line unless they are named. In \cref{sec:stepanov,sec:families}
we use the development's names for the main propositions (in sans-serif, e.g. $\sfname{VWin}$), so that
the statements can be compared with their formal counterparts.
